\section{Introduction}
\label{sec:introduction}

Ghana's external sector remains strongly exposed to commodity markets. UN Trade and Development classifies an economy as commodity-dependent when commodities account for more than 60\% of merchandise exports, and its 2025 country profile reports that Ghana's three leading commodity exports represented 78.4\% of allocated product exports over 2021--2023. Bank of Ghana data show the continuing importance of cocoa, gold and crude oil: in 2025, exports of cocoa beans and products were about US$4.0 billion, non-monetary gold about US$21.0 billion, and crude oil about US$2.6 billion. These exposures make the relation between commodity prices and the Ghanaian cedi relevant not only for macroeconomic analysis but also for financial risk measurement and decisions made under extreme market conditions \citep{unctad2025,bog2025}.

The Ghana-specific literature already shows that this relation is neither simple nor constant. \citet{archer2022} use monthly cocoa, gold, Brent crude oil and exchange-rate data and report asymmetric relationships across exchange-rate quantiles. \citet{barson2023} use wavelet methods and find that the comovement of Ghana's major commodity returns varies across time horizons, with the exchange rate playing a material role in those relationships. \citet{yeboah2025}, using quantile-on-quantile estimation over 2003--2023, likewise report nonlinear and asymmetric relationships among commodity prices, inflation and the exchange rate. Recent work has widened the setting further. \citet{atipaga2025} use daily data to examine the Cedi's connection with global uncertainty measures, while \citet{woode2025} study time- and frequency-varying connectedness among commodities, currencies and equity markets in commodity-dependent sub-Saharan African economies, including Ghana. Together, these studies make a fixed linear commodity-Cedi relationship difficult to justify.

International evidence points in the same direction. Copula studies of commodity currencies find that dependence can vary with the commodity, country, sample period and market state. \citet{reboredo2012} reports generally weak oil-exchange-rate dependence in the currencies studied, with stronger dependence after the global financial crisis. \citet{ignatieva2017} find positive dependence between commodity-price indices and the currencies of major commodity exporters, together with an increase in time-varying tail dependence after the crisis. More recent studies continue to find state dependence. \citet{go2024} document changing return and volatility spillovers between Malaysian crude palm oil and exchange rates, while \citet{dodd2026} show that the usual positive commodity-currency return relation can disappear when political risk is high. Crisis-focused work also provides reasons to examine the tails directly: \citet{qiao2023}, for example, report greater lower- and upper-tail contagion across commodity futures after the COVID-19 outbreak. Such findings motivate direct testing of stress-period tail behavior rather than assuming that dependence must increase in every crisis or market.

Copula methods are therefore well suited to the problem, but their use is not itself new in Ghana. \citet{mensah2020} estimate static and time-varying copulas for major foreign currencies against the Cedi and document tail-dependent exchange-rate comovement. The relevant gap is narrower. Existing Ghana studies have not centered the question of whether commodity-Cedi tail conclusions remain stable when the underlying daily data are aligned using different defensible market-calendar rules. That issue matters because commodity futures and the Cedi do not necessarily share the same trading days or observation timing. The classical nonsynchronous-trading literature shows that mismatched observation times can distort return-based dependence estimates. \citet{martens2001}, for example, show that close-to-close international stock returns can understate correlation and that synchronization changes estimated correlation dynamics. Related work finds that time adjustment can materially alter market-integration and spillover estimates. Calendar construction should therefore be treated as part of the statistical design rather than as an invisible data-cleaning step.

A second issue is model adequacy. Selecting the copula with the best relative information criterion does not establish that the selected family provides an adequate absolute representation of the data. Copula goodness-of-fit research distinguishes these questions and commonly relies on resampling when parameters are estimated. The present study therefore evaluates eight static bivariate copula families using a Rosenblatt-transform Cramér-von Mises statistic with parametric bootstrap, rather than treating relative fit as evidence of correctness. The marginal stage is also subject to explicit convergence and diagnostic gates because misspecified margins can distort copula and downstream risk estimates. This point is particularly relevant for the Cedi, whose return distribution contains many exact zeros and whose selected continuous marginal specification does not pass the same probability-integral-transform adequacy tests as the commodity margins. Cedi-related copula and extreme-value results are consequently interpreted as conditional diagnostics rather than model-free structural quantities.

A third issue concerns inference in the tails. The study uses empirical lower- and upper-tail concentration at prespecified finite thresholds rather than presenting these quantities as asymptotic tail-dependence coefficients. This distinction matters because finite-sample tail-dependence estimates are sensitive to the estimator and threshold. Stress-period comparisons are evaluated with a moving-block bootstrap, and multiplicity is controlled within each prespecified 60-test family using Bonferroni and Benjamini-Hochberg procedures. This design responds to a practical problem in large tail-comparison exercises: a collection of nominal p-values can look persuasive even when no result survives correction. Recent formal work on differences and equivalence between tail copulas, including \citet{bormann2020}, \citet{can2024}, and \citet{koike2025}, reinforces the need to state precisely which tail quantity is being compared and what inferential claim a test can support.

Against this background, the paper uses two co-equal constructions of the same daily-source data. The first preserves a business-day calendar and allows limited forward filling before returns are computed. The second computes close-to-close returns only between consecutive dates on which all five series have actual observations. Neither construction is treated as the truth or as a correction of the other. Instead, substantive findings are classified as calendar-robust when they survive both constructions, calendar-sensitive when the conclusions differ, and panel-specific when they arise only as a diagnostic of one construction. This makes calendar sensitivity an explicit part of the evidence rather than a post hoc robustness exercise.

The analysis addresses three questions. First, do empirical lower- and upper-tail concentration measures differ between calm periods and the COVID-19 and 2024 stress periods, and are those conclusions robust to alternative calendar-alignment methods? Second, can any of eight tested static bivariate copula families adequately represent the dependence structure among Ghana-relevant commodity and exchange-rate pairs, and which adequacy conclusions are robust across calendar constructions? Third, how sensitive are Cedi-related dependence and extreme-value diagnostics to the Cedi's zero-heavy return distribution, marginal-model inadequacy and calendar construction? By answering these questions jointly, the paper shifts attention from selecting a single preferred dependence model to identifying which conclusions remain defensible after calendar construction, marginal adequacy, absolute copula fit and multiple testing are taken seriously.
